% =========================================
% CHAPTER 2: BACKGROUND
% =========================================
\chapter{Background: Layer~2 Scaling Solutions}
\label{ch:background}

This chapter covers the technical concepts the rest of the thesis
builds on. Section~\ref{sec:fundamentals} describes the throughput
constraints of Layer~1 blockchains. Section~\ref{sec:l2concept}
introduces the general Layer~2 execution model. Sections~\ref{sec:zk}
and~\ref{sec:optimistic} describe each rollup family and note how this
thesis instantiates them. Section~\ref{sec:summary} connects the
chapter to the work that follows.

% -----------------------------------------
\section{Blockchain Fundamentals}
\label{sec:fundamentals}
% -----------------------------------------

Bitcoin~\cite{nakamoto2008bitcoin} established a decentralized
consensus model for exchanging value without intermediaries. Ethereum
extended it with a programmable execution layer~\cite{wood2014ethereum}.
Both networks face a hard throughput ceiling: Bitcoin sustains roughly
seven transactions per second, Ethereum around fifteen, both governed
by block-level resource limits~\cite{croman2016scaling}.

A blockchain can be modeled as a state machine. At each block~$t$ the
network holds a global state~$S_t$; incoming transactions produce a
successor:
\begin{equation}
  S_{t+1} = \text{apply}(S_t,\ \text{txs}_t)
\end{equation}
Integrity rests on three properties~\cite{wood2014ethereum}:
\textit{consensus} (agreement on transaction ordering),
\textit{finality} (irreversibility of confirmed blocks), and
\textit{data availability} (full auditability by any node). Layer~2
systems inherit these properties from the base layer rather than
reproducing them independently.

% -----------------------------------------
\section{The Layer~2 Execution Model}
\label{sec:l2concept}
% -----------------------------------------

Layer~2 systems separate execution from settlement~\cite{gudgeon2020sok}.
A \textit{sequencer} orders transactions, executes them off-chain, and
periodically submits a state commitment to a Layer~1 smart contract.
Security derives from that on-chain commitment, not from the off-chain
environment itself.

The economic motivation is cost amortization. One Layer~1 verification
is shared across an entire batch:
\begin{equation}
  \text{Cost per L2 tx} =
    \frac{\text{L1 batch cost} + \text{L2 execution cost}}{\text{batch size}}
\end{equation}
The two dominant approaches differ in what the sequencer submits
alongside the state root: a cryptographic validity proof (ZK~Rollups)
or nothing, relying on a post-hoc fraud challenge
(Optimistic~Rollups)~\cite{sok-layer2}.

% -----------------------------------------
\section{ZK Rollups}
\label{sec:zk}
% -----------------------------------------

A ZK~Rollup requires a validity proof with every state
commitment~\cite{groth2016size}. The prover computes a zk-SNARK
attesting that the transition from $\text{state}_\text{old}$ to
$\text{state}_\text{new}$ is correct:
\begin{equation}
  \pi = \text{Prove}(\text{state}_\text{old},\ \text{txs},\
    \text{state}_\text{new})
\end{equation}
The Layer~1 verifier contract accepts or rejects the submission
immediately based on the proof alone---no waiting period is
needed~\cite{groth2016size}. Production deployments such as
zkSync~\cite{zksync-paper} and Scroll~\cite{scroll2023} implement
Groth16~\cite{groth2016size} verification using Ethereum's BN254
pairing precompile.

\textbf{Security.}
Safety follows from proof soundness: a dishonest sequencer cannot
produce a valid proof for an incorrect state. This property is
sometimes called the \textit{one-honest-prover}
condition~\cite{chaliasos2024sok}---the system remains safe as long as
one correct prover exists. Finality is achieved shortly after the
proof lands on Layer~1, typically within one block confirmation
window.

\textbf{Constraints.}
Proof generation is computationally expensive, requiring specialized
hardware in production and adding measurable latency per
batch~\cite{zksync-paper}. On-chain verification costs between 200K
and 500K gas due to pairing operations~\cite{groth2016size}. Newer
schemes such as PLONK~\cite{gabizon2019plonk} and
STARKs~\cite{ben-sasson2018stark} reduce some of these overheads but
Groth16 remains the most widely deployed.

\textbf{Framework instantiation.}
\texttt{ZKVerifier.sol} accepts a state root and a mock proof,
binding it to a verifying-key hash. Proof generation is simulated
in \texttt{l2-zk/sequencer.py} using SHA-256 hashing. Batch
size and polling interval are configurable but default to 100~transactions
and 5~seconds, reflecting the higher per-batch submission cost relative to the Optimistic side.

% -----------------------------------------
\section{Optimistic Rollups}
\label{sec:optimistic}
% -----------------------------------------

An Optimistic~Rollup posts a state root without a proof and opens a
\textit{challenge window}---a period during which any observer can
dispute the commitment by submitting a fraud
proof~\cite{arbitrum-original}. In production deployments (Optimism,
Arbitrum) this window is seven days. In this framework it is
compressed to ten Layer~1 blocks ($\approx$2~minutes on a local
Hardhat node) for experimental tractability.

Fraud proving is implemented as a bisection protocol: the challenger
and sequencer narrow their disagreement to a single execution step,
which the Layer~1 contract re-executes to resolve the
dispute~\cite{arbitrum-original}. Optimism implements this via the
Cannon fault-proof VM~\cite{optimism2023cannon}; Arbitrum uses a
comparable WebAssembly-based mechanism~\cite{arbitrum-original}.

A commitment that survives its window without a valid challenge
transitions from optimistic acceptance to finality. This framework's
\texttt{OptimisticVerifier.sol} models the lifecycle through four
states: \texttt{PENDING}, \texttt{SOFT\_FINAL}, \texttt{HARD\_FINAL},
and \texttt{CHALLENGED}.

\textbf{Security.}
Safety depends on what Kalodner et al.~\cite{arbitrum-original} term
the \textit{at-least-one-honest-watcher} condition: the system is
secure as long as one participant monitors the chain and challenges
invalid roots within the window. This is an economic guarantee,
enforced through sequencer collateral, rather than a cryptographic
one~\cite{chaliasos2024sok}.

\textbf{Constraints.}
The challenge window is the primary usability cost: withdrawals to
Layer~1 are delayed by the full window, unless a liquidity provider
offers an early exit for a fee. Full transaction data must be published
on-chain so watchers can reconstruct state independently; EIP-4844
blob transactions~\cite{eip4844} reduce this data cost in recent
deployments.

\textbf{Framework instantiation.}
\texttt{l2-optimistic/sequencer.py} submits state roots without proof
generation. Batch size and polling interval are configurable but default to 200~transactions
and 8~seconds. Recovery scenarios are
triggered programmatically and observed through Prometheus metrics
(Chapter~\ref{ch:implementation}).

% -----------------------------------------
\section{Summary}
\label{sec:summary}
% -----------------------------------------

Both rollup families move execution off-chain and anchor security
on-chain, but differ in the mechanism that links the two: immediate
cryptographic verification in the ZK case versus deferred economic
enforcement in the Optimistic case. This single design choice drives
every dimension this thesis measures---finality time, recovery cost,
and security mechanism overhead. Chapter~\ref{ch:relatedwork} surveys
how existing work has studied these properties and identifies the gap
this framework addresses.